\documentclass[10pt]{article}
\usepackage[T1]{fontenc}
\usepackage[utf8]{inputenc}
\usepackage[french]{babel}
\usepackage{amsmath,amssymb,mathrsfs,mathpazo}
\usepackage{geometry}
\geometry{a4paper,margin=2cm}
\usepackage{graphicx,array,tabularx,tabularray,colortbl,multirow,multicol,booktabs,enumitem,pifont,ulem,fancybox,diagbox,cancel}
\usepackage{pstricks,pstricks-add,pst-plot,pst-tree,pst-node,pst-text,pst-func,pst-math,pst-bspline,pst-3dplot,pst-eucl,pst-solides3d}
\let\pscalculate\undefined
\usepackage{pst-bezier}
\usepackage{tikz,pgfplots,tkz-tab}
\usetikzlibrary{arrows,arrows.meta,patterns,calc,positioning,trees,decorations.markings,shapes.geometric,shapes.misc,shapes.arrows,intersections,decorations.pathmorphing,decorations.pathreplacing,backgrounds,shadows,babel}
\pgfplotsset{compat=1.18}
\usepgfplotslibrary{fillbetween}
\usepackage[np]{numprint}
\usepackage{listings,makecell,pas-tableur,twoopt}
\let\barre\undefined
\usepackage[pstricks]{tablvar}
\definecolor{darkgreen}{RGB}{0,100,0}
\providecommand{\red}{\color{red}}
\providecommand{\blue}{\color{blue}}
\providecommand{\green}{\color{green}}
\providecommand{\black}{\color{black}}
\providecommand{\white}{\color{white}}
\providecommand{\gray}{\color{gray}}
\providecommand{\cyan}{\color{cyan}}
\providecommand{\magenta}{\color{magenta}}
\providecommand{\yellow}{\color{yellow}}
\providecommand{\violet}{\color{violet}}
\providecommand{\orange}{\color{orange}}
\providecommand{\euro}{\texteuro}
\providecommand{\EUR}{\texteuro}
\providecommand{\ieme}{\textsuperscript{e}}
\providecommand{\degres}{\ensuremath{^\circ}}
\providecommand{\degree}{\ensuremath{^\circ}}
\providecommand{\vv}[1]{\overrightarrow{#1}}
\providecommand{\Coord}[1]{#1}
\providecommand{\decosix}{}
\providecommand{\starredbullet}{\ensuremath{\bullet}}
\pagestyle{empty}
\setlength{\parindent}{0pt}
\usepackage{adjustbox}\DeclareUnicodeCharacter{2236}{:}\DeclareUnicodeCharacter{2190}{\ensuremath{\leftarrow}}
\providecommand{\R}{\mathbb{R}}
\providecommand{\N}{\mathbb{N}}
\providecommand{\Z}{\mathbb{Z}}
\providecommand{\Q}{\mathbb{Q}}
\providecommand{\D}{\mathbb{D}}
\providecommand{\C}{\mathbb{C}}
\providecommand{\e}{{\mathrm{e}}}
\providecommand{\vect}[1]{\overrightarrow{#1}}
\providecommand{\vectt}[1]{\overrightarrow{\mathrm{#1}}}
\providecommand{\vvt}[1]{\overrightarrow{\mathrm{#1}}}
\providecommand{\np}[2][]{{#2~#1}}
\providecommand{\vv}[1]{{\overrightarrow{#1}}}
\renewcommand{\d}{{\mathrm{d}}}
\providecommand{\degres}{{^\circ}}
\providecommand{\degre}{{^\circ}}
\providecommand{\blue}{\color{blue}}
\providecommand{\red}{\color{red}}
\providecommand{\enskip}{\quad}
\providecommand{\llbracket}{\lbrack\!\lbrack}
\providecommand{\rrbracket}{\rbrack\!\rbrack}
\providecommand{\euro}{{\text{€}}}
\providecommand{\no}{{\text{n°}}}
\providecommand{\interval}[2]{{[#1;#2]}}
\providecommand{\intervFO}[2]{{[#1;#2[}}
\providecommand{\dotfill}{\ldots}

\providecommand{\ds}{\displaystyle}
\providecommand{\textding}[1]{#1}
\providecommand{\Oij}{\(\mathrm{O};\vec{\imath},\vec{\jmath}\)}
\providecommand{\Oijk}{\(\mathrm{O};\vec{\imath},\vec{\jmath},\vec{k}\)}
\providecommand{\Ouv}{\(\mathrm{O};\vec{u},\vec{v}\)}
\providecommand{\intervFF}[2]{\ensuremath{[#1;#2]}}
\renewcommand{\euro}{\ifmmode\text{\texteuro}\else\texteuro\fi}
\DeclareUnicodeCharacter{00B0}{\ifmmode^{\circ}\else\textdegree\fi}
\DeclareUnicodeCharacter{20AC}{\ifmmode\text{\texteuro}\else\texteuro\fi}
% Énoncés : Éducation nationale. Transcription et corrigés : APMEP, Denis Vergès et contributeurs. Découpage et classement : CoopMaths, Sébastien Lozano, Éric Elter et contributeurs. Adaptation de la présentation : Maths Explorer.
% CC BY-SA : https://coopmaths.fr/www/about/
% https://coopmaths.fr/alea/static/bac/2023/tex/bac_2023_09_sujet1_metropole_3.tex
\begin{document}
\noindent\textit{APMEP / CoopMaths. CC BY-SA. Présentation adaptée par Maths Explorer.}\par\medskip
\medskip

\emph{Les parties \textbf{A} et \textbf{B} sont indépendantes.\\
Les probabilités demandées seront données à $10^{-3}$ près.}

Pour aider à la détection de certaines allergies, on peut procéder à un test sanguin dont le résultat est soit positif, soit négatif.

Dans une population, ce test donne les résultats suivants :
\begin{itemize}
\item Si un individu est allergique, le test est positif dans 97\,\% des cas;
\item Si un individu n'est pas allergique, le test est négatif dans 95,7\,\% des cas.
\end{itemize}

Par ailleurs, 20\,\% des individus de la population concernée présentent un test positif.

On choisit au hasard un individu dans la population, et on note :

\begin{itemize}
\item[$\bullet~~$] $A$ l'évènement \og l'individu est allergique \fg{} ;
\item[$\bullet~~$] $T$ l'évènement \og l'individu présente un test positif \fg.
\end{itemize}

On notera $\overline{A}$ et $\overline{T}$ les évènements contraires de $A$ et $T$.
On appelle par ailleurs $x$ la probabilité de l'évènement $A$ : $x = p(A)$.

\begin{minipage}{9cm}
\textbf{Partie A}

\medskip

\begin{enumerate}
\item Reproduire et compléter l'arbre ci-contre décrivant la situation,
en indiquant sur chaque branche la probabilité correspondante.
\item 
	\begin{enumerate}
		\item Démontrer l'égalité : $p(T) = 0,927x + 0,043$.
		\item En déduire la probabilité que l'individu choisi soit allergique.
	\end{enumerate}	
\item Justifier par un calcul l'affirmation suivante :

\og Si le test d'un individu choisi au hasard est positif, il y a plus de 80\,\% de chances que cet individu soit allergique \fg.
\end{enumerate}
\end{minipage}\hfill
\begin{minipage}{5.5cm}
\begin{center}
\pstree[treemode=R,nodesepA=0pt,nodesepB=2.5pt,treesep = 1cm,levelsep=2.5cm]{\TR{}}
{\pstree{\TR{$A$~}\taput{$x$}}
{\TR{$T$}\taput{\ldots}
\TR{$\overline{T}$}\tbput{\ldots}
}
\pstree{\TR{$\overline{A}$~}\tbput{\ldots}}
{\TR{$T$}\taput{\ldots}
\TR{$\overline{T}$}\tbput{\ldots}
}
}
\end{center}
\end{minipage}

\textbf{Partie B}

\medskip

On réalise une enquête sur les allergies dans une ville en interrogeant $150$~habitants choisis au hasard, et on admet que ce choix se ramène à des tirages successifs indépendants avec remise.

On sait que la probabilité qu'un habitant choisi au hasard dans cette ville soit allergique est égale à $0,08$.

On note $X$ la variable aléatoire qui à un échantillon de $150$~habitants choisis au hasard associe le nombre de personnes allergiques dans cet échantillon.

\medskip

\begin{enumerate}
\item Quelle est la loi de probabilité suivie par la variable aléatoire $X$ ?

Préciser ses paramètres.
\item Déterminer la probabilité que $20$~personnes exactement parmi les $150$ interrogées soient allergiques.
\item Déterminer la probabilité qu'au moins 10\,\% des personnes parmi les $150$ interrogées soient allergiques.
\end{enumerate}

\bigskip

\end{document}
